% Generated draft: requires mathematical and source review.
\documentclass[11pt]{article}
\usepackage[margin=25mm]{geometry}
\usepackage{fontspec}
\setmainfont{Latin Modern Roman}
\usepackage{amsmath,amssymb,mathtools,unicode-math}
\setmathfont{Latin Modern Math}
\usepackage{enumitem,needspace,xurl,hyperref}
\setlength{\parindent}{0pt}
\setlength{\parskip}{5pt}
\setlength{\emergencystretch}{3em}
\newcommand{\N}{\mathbb{N}}
\newcommand{\Z}{\mathbb{Z}}
\newcommand{\Q}{\mathbb{Q}}
\newcommand{\R}{\mathbb{R}}
\newcommand{\C}{\mathbb{C}}
\newcommand{\F}{\mathbb{F}}
\newcommand{\A}{\mathbb{A}}
\newcommand{\PP}{\mathbb{P}}
\newcommand{\E}{\mathbb{E}}
\newcommand{\eps}{\varepsilon}
\newcommand{\dd}{\,\mathrm{d}}
\begin{document}
\section*{UnsolvedMath ID: 30006392}
\textbf{Source problem code:} catalogue entry
\subsection{Definitions and mathematical statement}
Let \(G\) be a finite group with identity \(e=\mathrm{id}\). For a subset \(S\subseteq G\setminus\{e\}\), put \(m=|S|\). An ordering of \(S\) is a sequence \((s_1,\ldots,s_m)\) containing each element of \(S\) exactly once. Define its nonempty successive partial products, in the displayed order, by
\[
p_1=s_1,\qquad p_k=p_{k-1}s_k=s_1s_2\cdots s_k\quad(2\leq k\leq m).
\]
No commutativity is assumed. Call the ordering valid, under the convention used here, if
\[
p_i\neq p_j\qquad\text{for all }1\leq i<j\leq m.
\]
For \(S=\varnothing\), the empty ordering satisfies this condition vacuously.

\textbf{Reported question.} Characterize the finite groups \(G\) satisfying
\[
\forall S\subseteq G\setminus\{e\},\quad
\exists\text{ an ordering }(s_1,\ldots,s_{|S|})\text{ of }S
\quad\text{such that}\quad
s_1\cdots s_i\neq s_1\cdots s_j
\quad(1\leq i<j\leq |S|).
\]
Does every finite group satisfy this property?

\textbf{Unresolved source qualifications.} This formulation interprets ``successive partial products'' as the nonempty products \(p_1,\ldots,p_m\). The supplied excerpt does not reproduce the definition of ``valid ordering,'' so it does not establish whether the empty partial product \(p_0=e\) was also intended to participate in the distinctness condition. Including it would additionally require \(p_k\neq e\) for every \(1\leq k\leq m\), and would give a different formulation. The displayed question above therefore remains conditional on the nonempty-product convention. The original excerpt also contains a detached footnote fragment referring to an ordering with total product equal to \(\mathrm{id}\); its antecedent is missing. No total-product hypothesis is imposed here, because its intended role cannot be recovered from the supplied text.

The registry labels this problem open and its edited statement \texttt{corrected\_verified}. These are supplied curation labels, not independent mathematical or source verification by this draft.
\subsection{Short English statement}
Determine which finite groups have the following property: every subset excluding the identity can be listed so that its nonempty successive partial products are pairwise distinct. The supplied question asks whether all finite groups have this property. The intended treatment of the empty partial product, and the role of a detached total-product footnote, require source review.
\subsection{Sources}
\begin{itemize}
\item Supplied local UnsolvedMath registry record, permanent ID \(30006392\), title ``Valid Orderings of Subsets of Finite Groups,'' source problem code \texttt{catalogue entry}, in the Oberwolfach Reports Open Problems collection. The record gives proposed year \(2025\), no named proposer, status \texttt{open}, and statement status \texttt{corrected\_verified}. Its verification note reports that the official report was checked and the question isolated and its notation repaired; this draft has not independently checked that report.
\item Supplied report citation: \emph{Combinatorics, Probability and Computing (2026)}, Oberwolfach Reports, DOI \texttt{10.4171/owr/2025/42}, pp.~2243--2322. \url{https://doi.org/10.4171/owr/2025/42}. The supplied background describes these reports as workshop records including open-problem discussions. The citation's \(2026\) designation, the DOI's \(2025\) component, and the registry's proposed year \(2025\) are retained as supplied; their relationship has not been independently verified.
\item Attribution appearing in the supplied original excerpt: ``On Graham's rearrangement conjecture over \(\F_2^n\),'' Matija Buci\'{c}, joint work with Benjamin Bedert, Noah Kravitz, Richard Montgomery, and Alp M\"uyesser. This is attribution from the excerpt, not a verified identification of the problem's proposer.
\item Supplied literature source: \emph{On Graham's Rearrangement Conjecture over \(\F_2^n\)}. \url{https://arxiv.org/abs/2508.18254}.
\item Supplied literature source: \emph{Towards Graham's Rearrangement Conjecture via Rainbow Paths}. \url{https://www.sciencedirect.com/science/article/pii/S0001870826001143}.
\item Literature assessment dated exactly \texttt{2026-08-22}: the supplied curation labels the general finite-group subset-ordering problem unresolved and reports that work over \(\F_2^n\) establishes valid orderings for all sufficiently large subsets without settling arbitrary subsets in arbitrary finite groups. This is supplied dated curation derived from the OWR material, not an independently verified assessment. The record itself requests independent verification of time-sensitive status and cited claims; no precise threshold or quantified version of the reported partial result is supplied.
\end{itemize}
\end{document}
